\documentclass[11pt,a4paper]{article}

\usepackage[T1]{fontenc}
\usepackage[utf8]{inputenc}
\usepackage{lmodern}
\usepackage[a4paper,margin=2.2cm]{geometry}
\usepackage{setspace}
\usepackage{enumitem}
\usepackage[hidelinks]{hyperref}
\usepackage{titlesec}

\setstretch{1.05}
\setlist[itemize]{topsep=3pt,itemsep=2pt,parsep=0pt}
\setlist[enumerate]{topsep=3pt,itemsep=2pt,parsep=0pt}

\titleformat{\section}{\large\bfseries}{\thesection.}{0.6em}{}
\titleformat{\subsection}{\normalsize\bfseries}{\thesubsection}{0.6em}{}
\titleformat{\subsubsection}{\normalsize\itshape}{\thesubsubsection}{0.6em}{}

\title{Chapter 18 --- Central Nervous System\\Structured Rewrite}
\author{}
\date{}

\begin{document}

\maketitle

\noindent\textbf{Source:} Peck \& Harris, \textit{Pharmacology for Anaesthesia and Intensive Care}, 5th ed., Chapter 18: Central Nervous System, pp.\ 253--265.

\tableofcontents

\newpage

\section{Hypnotics and anxiolytics}

\subsection{Physiology}

\subsubsection{GABA}
\begin{itemize}
  \item GABA is the main inhibitory neurotransmitter in the CNS.
  \item It acts via two receptor subtypes:
\end{itemize}

\paragraph{GABA\textsubscript{A}}
\begin{itemize}
  \item Ligand-gated chloride ion channel.
  \item Five subunits arranged around a central ion channel.
  \item GABA increases opening frequency of the chloride channel.
  \item Result: increased chloride conductance and neuronal hyperpolarisation.
  \item Benzodiazepines potentiate chloride conductance by binding to the $\alpha$ subunit of the activated receptor complex.
  \item Mainly postsynaptic and widely distributed through the CNS.
\end{itemize}

\paragraph{GABA\textsubscript{B}}
\begin{itemize}
  \item Metabotropic receptor acting via G-protein / second messenger systems.
  \item Increases potassium conductance.
  \item Causes neuronal hyperpolarisation.
  \item Present both presynaptically and postsynaptically.
  \item Also present in the dorsal horn of the spinal cord.
  \item Baclofen acts via GABA\textsubscript{B} receptors to reduce spasticity.
\end{itemize}

\subsubsection{Benzodiazepine receptor subtype concept}
\begin{itemize}
  \item Benzodiazepines modulate GABA at GABA\textsubscript{A} receptors.
  \item The chapter links effect to the $\alpha$-subunit type.
  \item Two receptor subtypes are described:
  \begin{itemize}
    \item \textbf{BZ1}: spinal cord and cerebellum; responsible for anxiolysis.
    \item \textbf{BZ2}: spinal cord, hippocampus and cerebral cortex; responsible for sedative and anticonvulsant activity.
  \end{itemize}
\end{itemize}

\subsection{Benzodiazepines}

\subsubsection{Uses of benzodiazepines as a class}
\begin{itemize}
  \item Premedication in anaesthesia.
  \item Sedation for minor procedures.
  \item More widely as anxiolytics, hypnotics, and anticonvulsants.
\end{itemize}

\subsubsection{Structure}
\begin{itemize}
  \item Core structure:
  \begin{itemize}
    \item one benzene ring
    \item one seven-membered di-azepine ring
  \end{itemize}
  \item For pharmacological activity they also require:
  \begin{itemize}
    \item a carbonyl group at position 2
    \item another benzene ring
    \item a halogen on the first benzene ring
  \end{itemize}
\end{itemize}

\subsection{Midazolam}

\subsubsection{Presentation}
\begin{itemize}
  \item Clear solution at pH 3.5.
  \item Unique among benzodiazepines because its structure depends on surrounding pH.
  \item At pH 3.5:
  \begin{itemize}
    \item di-azepine ring is open
    \item molecule is ionised
    \item therefore water-soluble
  \end{itemize}
  \item At pH $>$ 4:
  \begin{itemize}
    \item ring closes
    \item molecule becomes non-ionised
    \item therefore lipid-soluble
  \end{itemize}
  \item pKa 6.5.
  \item At physiological pH, 89\% is unionised.
  \item Because it is water-soluble in formulation, it does not cause pain on injection.
\end{itemize}

\subsubsection{Uses}
\begin{itemize}
  \item Intravenous sedation for minor procedures.
  \item Has powerful amnesic properties.
  \item Used to sedate ventilated ICU patients.
\end{itemize}

\subsubsection{Kinetics}
\begin{itemize}
  \item Can be given orally, intranasally, or intramuscularly.
  \item Oral bioavailability about 40\%.
  \item Short duration of action due to redistribution.
  \item Metabolised by hydroxylation to 1-$\alpha$-hydroxymidazolam (active).
  \item Then conjugated with glucuronic acid before renal excretion.
  \item Less than 5\% becomes oxazepam.
  \item Protein binding about 95\%.
  \item Elimination half-life 1--4 h.
  \item Clearance 6--10 mL$\cdot$kg$^{-1}\cdot$min$^{-1}$.
  \item Compared with diazepam and lorazepam:
  \begin{itemize}
    \item higher clearance
    \item therefore effects wear off more rapidly after infusion
  \end{itemize}
  \item Interaction highlighted:
  \begin{itemize}
    \item alfentanil is metabolised by the same hepatic CYP isoenzyme (3A3/4)
    \item co-administration may prolong effects
  \end{itemize}
\end{itemize}

\subsection{Diazepam}

\subsubsection{Key points emphasised by the chapter}
\begin{itemize}
  \item High lipid solubility:
  \begin{itemize}
    \item facilitates oral absorption
    \item gives rapid central effects
  \end{itemize}
  \item Highly protein-bound ($\sim$95\% to albumin).
  \item Hepatic oxidation to:
  \begin{itemize}
    \item desmethyldiazepam
    \item oxazepam
    \item temazepam
  \end{itemize}
  \item All are active metabolites.
  \item Does not induce hepatic enzymes.
  \item Glucuronide derivatives excreted in urine.
  \item Of the benzodiazepines compared here, it has the lowest clearance.
  \item Its half-life is hugely increased by infusion.
\end{itemize}

\subsubsection{Clinical / adverse effect points}
\begin{itemize}
  \item May cause some cardiorespiratory depression.
  \item Liver failure and cimetidine prolong its action by reducing metabolism.
  \item With opioids or alcohol, respiratory depression may be more marked.
  \item Like other benzodiazepines, it reduces MAC of co-administered anaesthetic agents.
\end{itemize}

\subsection{Lorazepam}

\subsubsection{What the chapter stresses}
\begin{itemize}
  \item Similar pharmacokinetics and actions to other benzodiazepines.
  \item Important distinction: its metabolites are inactive.
\end{itemize}

\subsubsection{Uses}
\begin{itemize}
  \item Premedication.
  \item Anxiolytic.
  \item Amnesic.
  \item First-line therapy for status epilepticus.
\end{itemize}

\subsubsection{Kinetics}
\begin{itemize}
  \item Well absorbed orally and intramuscularly.
  \item Highly protein-bound ($\sim$95\%).
  \item Conjugated with glucuronic acid.
  \item Produces inactive metabolites.
  \item Excreted in urine.
\end{itemize}

\subsection{Temazepam}

\subsubsection{Uses}
\begin{itemize}
  \item Night-time sedative.
  \item Anxiolytic premedicant.
\end{itemize}

\subsubsection{Chapter emphasis}
\begin{itemize}
  \item ``No unique features'' within the benzodiazepine family.
\end{itemize}

\subsubsection{Kinetics}
\begin{itemize}
  \item Well absorbed from the gut.
  \item 75\% protein-bound.
  \item Volume of distribution 0.8 L$\cdot$kg$^{-1}$.
  \item 80\% excreted unchanged in urine.
  \item Hepatic glucuronidation occurs.
  \item Only a very small amount is demethylated to oxazepam.
  \item Half-life about 8 h.
  \item May cause hangover effects.
\end{itemize}

\subsection{Chlordiazepoxide}

\subsubsection{What the chapter attaches specifically to this drug}
\begin{itemize}
  \item First benzodiazepine derivative available for clinical use.
  \item Anxiolytic sedative.
  \item Used to treat symptoms associated with alcohol withdrawal.
\end{itemize}

\subsubsection{Dose note given in the text}
\begin{itemize}
  \item For short-term anxiety:
  \begin{itemize}
    \item 5 mg tds
    \item up to 25 mg qds
  \end{itemize}
  \item Should be titrated carefully to individual need.
\end{itemize}

\subsubsection{Kinetics}
\begin{itemize}
  \item High oral bioavailability.
  \item Complex metabolic pathway with multiple active metabolites.
  \item Elimination half-life 5--30 h.
  \item Half-life may be increased in the elderly.
\end{itemize}

\subsection{Flumazenil}

\subsubsection{Classification}
\begin{itemize}
  \item Imidazobenzodiazepine.
\end{itemize}

\subsubsection{Uses}
\begin{itemize}
  \item Reversal of benzodiazepine effects:
  \begin{itemize}
    \item excessive sedation after minor procedures
    \item benzodiazepine overdose
  \end{itemize}
  \item Important caution:
  \begin{itemize}
    \item in mixed drug overdose it may precipitate fits
  \end{itemize}
  \item Given intravenously in 100 microgram increments.
  \item Acts within 2 minutes.
  \item Half-life about 1 h.
  \item Because it is shorter acting than many benzodiazepines, repeat doses or infusion may be required.
\end{itemize}

\subsubsection{Mechanism of action}
\begin{itemize}
  \item Competitive benzodiazepine antagonist.
  \item The text adds nuance:
  \begin{itemize}
    \item has some agonist activity
    \item ability to precipitate seizures in certain patients may reflect inverse agonist activity
  \end{itemize}
\end{itemize}

\subsubsection{Effects}
\begin{itemize}
  \item Nausea and vomiting.
  \item May precipitate:
  \begin{itemize}
    \item anxiety
    \item agitation
    \item seizures
  \end{itemize}
  \item Especially in epileptic patients.
\end{itemize}

\subsubsection{Kinetics}
\begin{itemize}
  \item 50\% protein-bound.
  \item Significant hepatic metabolism.
  \item Metabolites are inactive.
  \item Excreted in urine.
\end{itemize}

\subsection{$\alpha$-agonists}

\subsubsection{Dexmedetomidine}

\paragraph{Presentation and uses}
\begin{itemize}
  \item S-enantiomer of medetomidine.
  \item R-form has no pharmacological activity.
  \item Used for light to moderate sedation.
  \item Especially for usually ventilated ICU patients.
  \item Also used for:
  \begin{itemize}
    \item procedural sedation, including paediatric patients
    \item balanced general anaesthesia
  \end{itemize}
  \item Presented as a colourless solution.
  \item Vials contain 100 microgram$\cdot$mL$^{-1}$.
  \item Typically diluted to 4--8 microgram$\cdot$mL$^{-1}$.
\end{itemize}

\paragraph{Dose}
\begin{itemize}
  \item ICU sedation:
  \begin{itemize}
    \item initial rate 0.7 microgram$\cdot$kg$^{-1}\cdot$h$^{-1}$
    \item adjusted to effect
    \item range 0.2--1.4 microgram$\cdot$kg$^{-1}\cdot$h$^{-1}$
  \end{itemize}
  \item In anaesthetic practice:
  \begin{itemize}
    \item loading dose 1 microgram$\cdot$kg$^{-1}$ over 10 min
    \item maintenance 0.2--0.6 microgram$\cdot$kg$^{-1}\cdot$h$^{-1}$
  \end{itemize}
  \item Much larger infusion rates (5--10 microgram$\cdot$kg$^{-1}\cdot$h$^{-1}$) described for rare use as a sole general anaesthetic.
\end{itemize}

\paragraph{Mechanism of action}
\begin{itemize}
  \item Highly selective $\alpha_2$-receptor agonist.
  \item $\alpha_2:\alpha_1$ ratio = 1600:1.
  \item Said to be eight times more selective than clonidine.
  \item Inhibits noradrenaline release at the locus coeruleus in the pons.
  \item The chapter links this to natural sleep pathways.
\end{itemize}

\paragraph{Effects}
\begin{itemize}
  \item \textbf{CNS}
  \begin{itemize}
    \item sedation
    \item anxiolysis
    \item sedation resembles physiological sleep more than pharmacological sedation/anaesthesia
    \item supported by EEG studies
    \item associated with less delirium than some drug protocols
    \item also produces analgesia via actions in the substantia gelatinosa of the dorsal horn, reducing release of transmitters involved in nociception
  \end{itemize}
  \item \textbf{Cardiovascular}
  \begin{itemize}
    \item biphasic
    \item initial hypertension via $\alpha_{2B}$ receptors on vascular smooth muscle
    \item then hypotension and bradycardia from reduced central noradrenaline release
    \item this pattern is more common during the loading dose
    \item rare progression to asystole reported when used as a sole agent
    \item should be used with extreme caution in conduction defects
  \end{itemize}
  \item \textbf{Respiratory}
  \begin{itemize}
    \item appears to cause no respiratory impairment
    \item airway remains patent
    \item may be particularly useful during extubation
    \item also has decongestant and antisialagogue effects
    \item useful in sedation for fibreoptic intubation
  \end{itemize}
\end{itemize}

\paragraph{Kinetics}
\begin{itemize}
  \item Oral absorption unpredictable:
  \begin{itemize}
    \item bioavailability $<$20\%
  \end{itemize}
  \item Buccal absorption better:
  \begin{itemize}
    \item bioavailability $>$75\%
  \end{itemize}
  \item These non-parenteral routes avoid the initial hypertension seen with parenteral use.
  \item Protein binding 94\%.
  \item Crosses easily into the brain.
  \item Undergoes both:
  \begin{itemize}
    \item Phase I hydroxylation
    \item Phase II glucuronidation and methylation
  \end{itemize}
  \item Inactive metabolites excreted in urine.
  \item Therefore hepatic impairment, not renal impairment, prompts dose reduction.
  \item Distribution half-life 6 min.
  \item Elimination half-life about 2 h.
  \item Clearance about 40 L$\cdot$h$^{-1}$.
  \item Volume of distribution 118 L.
  \item Vd increased in hypoalbuminaemia.
\end{itemize}

\subsection{Non-benzodiazepine hypnotics (Z drugs)}

\subsubsection{Drugs listed}
\begin{itemize}
  \item Zopiclone
  \item Zaleplon
  \item Zolpidem
\end{itemize}

\subsubsection{What the chapter emphasises}
\begin{itemize}
  \item Act via the benzodiazepine receptor.
  \item Not structurally benzodiazepines.
  \item Developed to try to reduce dependence and next-day sedation.
\end{itemize}

\subsubsection{Half-lives}
\begin{itemize}
  \item Zopiclone: $\sim$5 h
  \item Zaleplon: $\sim$1 h
  \item Zolpidem: $\sim$2 h
\end{itemize}

\subsubsection{Overall judgement given in the chapter}
\begin{itemize}
  \item They have not had a major impact.
  \item Carry essentially the same side effects as benzodiazepines.
  \item Switching is recommended only where there is intolerance to a specific benzodiazepine.
  \item With general intolerance, patients are likely to be intolerant of Z drugs as well.
\end{itemize}

\section{Antidepressants}

\subsection{Main groups}
\begin{itemize}
  \item Tricyclics
  \item Selective serotonin reuptake inhibitors
  \item Monoamine oxidase inhibitors
  \item Atypical agents
\end{itemize}

\subsection{Tricyclics}

\subsubsection{Drugs listed}
\begin{itemize}
  \item Amitriptyline
  \item Nortriptyline
  \item Imipramine
  \item Dothiepin
\end{itemize}

\subsubsection{Uses}
\begin{itemize}
  \item Depressive illness.
  \item Nocturnal enuresis.
  \item Adjunct in chronic pain treatment.
\end{itemize}

\subsubsection{Mechanism of action}
\begin{itemize}
  \item Competitive block of neuronal uptake 1 of:
  \begin{itemize}
    \item noradrenaline
    \item serotonin
  \end{itemize}
  \item Therefore increase synaptic transmitter concentration.
  \item Antidepressant effect takes up to 2 weeks.
  \item Also block:
  \begin{itemize}
    \item muscarinic receptors
    \item histaminergic receptors
    \item $\alpha$-adrenoceptors
  \end{itemize}
  \item Also have non-specific sedative effects.
\end{itemize}

\subsubsection{Effects}
\begin{itemize}
  \item \textbf{CNS}
  \begin{itemize}
    \item sedation
    \item occasionally seizures in epileptic patients
  \end{itemize}
  \item \textbf{Anticholinergic}
  \begin{itemize}
    \item dry mouth
    \item constipation
    \item urinary retention
    \item blurred vision
  \end{itemize}
  \item \textbf{Cardiovascular}
  \begin{itemize}
    \item postural hypotension, especially in the elderly
  \end{itemize}
\end{itemize}

\subsubsection{Kinetics}
\begin{itemize}
  \item Well absorbed from the gut.
  \item High lipid solubility.
  \item Highly protein-bound.
  \item High volume of distribution.
  \item Hepatic metabolism with marked interpatient variability.
  \item Often produce active metabolites:
  \begin{itemize}
    \item example given: imipramine $\rightarrow$ desipramine and nortriptyline
  \end{itemize}
\end{itemize}

\subsection{Tricyclic overdose}

\subsubsection{Features}
\begin{itemize}
  \item \textbf{Cardiovascular}
  \begin{itemize}
    \item sinus tachycardia common
    \item dose-related QT prolongation
    \item QRS widening
    \item ventricular arrhythmias more likely when QRS $>$ 0.16 s
    \item right bundle branch block may be seen
    \item blood pressure may be high or low
    \item severe overdose may give refractory hypotension and pulseless electrical activity
  \end{itemize}
  \item \textbf{Central}
  \begin{itemize}
    \item excitation
    \item seizures
    \item then depression
    \item seizures correlate with QRS $>$ 0.1 s
    \item mydriasis and hyperthermia also occur
  \end{itemize}
  \item Anticholinergic effects.
\end{itemize}

\subsubsection{Treatment}
\begin{itemize}
  \item Gastric lavage followed by activated charcoal.
  \item Supportive care.
  \item Seizures:
  \begin{itemize}
    \item benzodiazepines or phenytoin
  \end{itemize}
  \item Ventricular arrhythmias:
  \begin{itemize}
    \item phenytoin or lidocaine
  \end{itemize}
  \item Inotropes should be avoided where possible because of arrhythmia risk.
  \item Intravascular volume expansion is often sufficient for hypotension.
  \item Reversal of anticholinergic effects with an anticholinesterase is not recommended:
  \begin{itemize}
    \item may precipitate seizures, bradycardia and heart failure
  \end{itemize}
\end{itemize}

\subsection{SSRIs}

\subsubsection{Drugs listed}
\begin{itemize}
  \item Fluoxetine
  \item Paroxetine
  \item Sertraline
  \item Venlafaxine
\end{itemize}

\subsubsection{Class points}
\begin{itemize}
  \item Selectively inhibit neuronal reuptake of 5-HT.
  \item No more effective than standard antidepressants.
  \item But do not have the same side-effect profile.
  \item Compared with TCAs they are:
  \begin{itemize}
    \item less sedative
    \item less anticholinergic
    \item less cardiotoxic in overdose
  \end{itemize}
  \item Still associated with GI effects:
  \begin{itemize}
    \item nausea
    \item constipation
  \end{itemize}
  \item They may impair platelet function by inhibiting 5-HT uptake into platelets.
  \item This may matter in the presence of other antiplatelet agents.
\end{itemize}

\subsubsection{Serotonergic syndrome}
\begin{itemize}
  \item Potentially fatal.
  \item Characterised by:
  \begin{itemize}
    \item hyper-reflexia
    \item agitation
    \item clonus
    \item hyperthermia
  \end{itemize}
  \item Commonest combination:
  \begin{itemize}
    \item MAOI + SSRI
  \end{itemize}
  \item But phenylpiperidine opioids, especially pethidine, may also precipitate it because they have weak serotonin reuptake inhibitor properties.
\end{itemize}

\subsubsection{Fluoxetine}
\begin{itemize}
  \item Effective antidepressant with minimal sedation.
  \item 50:50 mix of two equally active isomers.
  \item Well absorbed.
  \item Metabolised in liver by CYP450 enzymes.
  \item Also has non-saturable enzymes that prevent unchecked rise in level.
  \item Dose should be reduced in renal failure because accumulation may occur.
  \item Side effects:
  \begin{itemize}
    \item nausea and vomiting
    \item headache
    \item insomnia
    \item reduced libido
    \item mania or hypomania in up to 1\%
  \end{itemize}
\end{itemize}

\subsubsection{Venlafaxine}
\begin{itemize}
  \item Appears to block reuptake of:
  \begin{itemize}
    \item noradrenaline
    \item 5-HT
    \item to a lesser extent dopamine
  \end{itemize}
  \item Has little effect on:
  \begin{itemize}
    \item muscarinic receptors
    \item histaminergic receptors
    \item $\alpha$-adrenoceptors
  \end{itemize}
\end{itemize}

\subsection{Monoamine oxidase inhibitors}

\subsubsection{Uses of the group}
\begin{itemize}
  \item Resistant depression
  \item Obsessive compulsive disorders
  \item Chronic pain syndromes
  \item Migraine
\end{itemize}

\subsubsection{MAO physiology in the chapter}
\begin{itemize}
  \item MAO is present as isoenzymes in presynaptic neurones.
  \item Responsible for deamination of amine neurotransmitters.
  \item Two types:
  \begin{itemize}
    \item \textbf{MAO-A}: prefers 5-HT and catecholamines
    \item \textbf{MAO-B}: prefers tyramine and phenylethylamine
  \end{itemize}
  \item Two generations of MAOI:
  \begin{itemize}
    \item old: irreversible, non-selective
    \item newer: selective, reversible MAO-A inhibitors (RIMA)
  \end{itemize}
  \item Not first-line because of serious side effects and hepatic toxicity.
\end{itemize}

\subsection{Non-selective irreversible MAOIs}

\subsubsection{Drugs listed}
\begin{itemize}
  \item Phenelzine
  \item Isocarboxazid
  \item Tranylcypromine
\end{itemize}

\subsubsection{Drug-specific chapter points}
\begin{itemize}
  \item Phenelzine and isocarboxazid are hydrazines.
  \item Tranylcypromine is a non-hydrazine.
  \item Tranylcypromine is potentially the most dangerous because it has stimulant activity.
\end{itemize}

\subsubsection{Effects}
\begin{itemize}
  \item Treat depression.
  \item Also produce:
  \begin{itemize}
    \item sedation
    \item blurred vision
    \item orthostatic hypotension
  \end{itemize}
  \item May cause hypertensive crises after:
  \begin{itemize}
    \item tyramine-rich foods
    \item indirectly acting sympathomimetics
  \end{itemize}
  \item Foods specifically listed:
  \begin{itemize}
    \item cheese
    \item pickled herring
    \item chicken liver
    \item Bovril
    \item chocolate
  \end{itemize}
  \item Hepatic enzymes are inhibited.
  \item Hydrazine compounds may cause hepatotoxicity.
  \item Interaction with pethidine may cause:
  \begin{itemize}
    \item cerebral irritability
    \item hyperpyrexia
    \item cardiovascular instability
  \end{itemize}
  \item Interaction with fentanyl has also been reported.
\end{itemize}

\subsection{Selective reversible MAOI-A}

\subsubsection{Moclobemide}
\begin{itemize}
  \item Causes less potentiation of tyramine than older MAOIs.
  \item Usually does not require the same degree of dietary restriction.
  \item However some patients are especially tyramine-sensitive.
  \item So the chapter still advises avoidance of:
  \begin{itemize}
    \item tyramine-rich foods
    \item indirectly acting sympathomimetic amines
  \end{itemize}
\end{itemize}

\subsubsection{Kinetics}
\begin{itemize}
  \item Completely absorbed from gut.
  \item Significant first-pass metabolism.
  \item Oral bioavailability 60--80\%.
  \item Metabolised in liver by CYP450.
  \item Slow metabolisers:
  \begin{itemize}
    \item about 2\% of Caucasian population
    \item about 15\% of Asian population
  \end{itemize}
  \item Metabolites excreted in urine.
\end{itemize}

\subsubsection{Linezolid}
\begin{itemize}
  \item Not an antidepressant, but the chapter explicitly notes it here.
  \item Used for MRSA and VRE.
  \item Is also a MAOI.
  \item Therefore carries the usual MAOI cautions and contraindications.
\end{itemize}

\subsection{MAOIs and general anaesthesia}

\subsubsection{Emergency surgery}
\begin{itemize}
  \item Do not give:
  \begin{itemize}
    \item pethidine
    \item indirectly acting sympathomimetics such as ephedrine
  \end{itemize}
  \item If cardiovascular support is needed:
  \begin{itemize}
    \item use direct-acting agents
    \item but with extreme caution
  \end{itemize}
\end{itemize}

\subsubsection{Elective surgery}
\begin{itemize}
  \item Stopping MAOIs 14--21 days before surgery risks relapse of depression.
  \item The chapter notes that newer agents may:
  \begin{itemize}
    \item control depression more effectively
    \item reduce the chance of serious peri-operative interaction
  \end{itemize}
\end{itemize}

\subsubsection{Interaction logic stressed by the text}
\begin{itemize}
  \item Indirectly acting sympathomimetics rely heavily on MAO for metabolism.
  \item Therefore can cause exaggerated:
  \begin{itemize}
    \item hypertension
    \item arrhythmias
  \end{itemize}
  \item Directly acting sympathomimetics should also be used cautiously.
  \item But they are also metabolised by COMT, so exaggerated response is less marked.
\end{itemize}

\subsubsection{Switching rules given}
\begin{itemize}
  \item Stop MAOIs for 2 weeks before alternative antidepressant therapy.
  \item Allow 2 weeks from the end of TCA therapy before starting MAOI therapy.
\end{itemize}

\subsection{Atypical agents}

\subsubsection{Mianserin}
\begin{itemize}
  \item Tetracyclic compound.
  \item Used in depressive illness, especially where sedation is required.
  \item Does not block neuronal reuptake of transmitters like TCAs.
  \item Instead blocks presynaptic $\alpha_2$-adrenoceptors.
  \item This reduces negative feedback and increases synaptic neurotransmitter concentration.
  \item Has very little ability to block:
  \begin{itemize}
    \item muscarinic receptors
    \item peripheral $\alpha$-adrenoceptors
  \end{itemize}
  \item Therefore causes less:
  \begin{itemize}
    \item antimuscarinic effect
    \item postural hypotension
  \end{itemize}
  \item Important adverse effects:
  \begin{itemize}
    \item agranulocytosis
    \item aplastic anaemia
  \end{itemize}
  \item These are more common in the elderly.
\end{itemize}

\section{Lithium carbonate}

\subsection{Uses}
\begin{itemize}
  \item Bipolar depression.
  \item Mania.
  \item Recurrent affective disorders.
\end{itemize}

\subsection{Key pharmacology}
\begin{itemize}
  \item Narrow therapeutic index.
  \item Plasma level should be 0.5--1.5 mmol$\cdot$L$^{-1}$.
  \item In excitable cells, lithium imitates sodium.
  \item Decreases release of neurotransmitters.
\end{itemize}

\subsection{Effects and adverse effects}
\begin{itemize}
  \item May increase generalised muscle tone.
  \item May lower seizure threshold in epileptics.
  \item Polyuria and polydipsia due to ADH antagonism.
  \item May raise serum:
  \begin{itemize}
    \item sodium
    \item magnesium
    \item calcium
  \end{itemize}
  \item Thyroid may enlarge and become underactive.
  \item Weight gain and tremor may occur.
\end{itemize}

\subsection{Anaesthetic relevance}
\begin{itemize}
  \item Prolongs neuromuscular blockade.
  \item May decrease anaesthetic requirements.
  \item Chapter explanation:
  \begin{itemize}
    \item blocks brainstem release of noradrenaline and dopamine
  \end{itemize}
\end{itemize}

\subsection{Toxicity above therapeutic level}
\begin{itemize}
  \item Vomiting
  \item Abdominal pain
  \item Ataxia
  \item Convulsions
  \item Arrhythmias
  \item Death
\end{itemize}

\section{Anticonvulsants}

\subsection{General introductory point}
\begin{itemize}
  \item Prefer single-agent therapy where possible.
  \item Avoids drug interactions.
  \item Patients rarely improve with a second agent.
\end{itemize}

\subsection{Phenytoin}

\subsubsection{Uses}
\begin{itemize}
  \item Grand mal seizures.
  \item Partial seizures.
  \item Trigeminal neuralgia.
  \item Ventricular arrhythmias after TCA overdose.
  \item Can be given orally or intravenously.
  \item Dose must be tailored because of wide interpatient variation.
  \item About 9\% of the population are slow hydroxylators.
  \item Blood assays are useful.
  \item Incompatible with 5\% dextrose because it becomes gelatinous.
  \item Therapeutic level 10--20 microgram$\cdot$mL$^{-1}$.
\end{itemize}

\subsubsection{Mechanism of action}
\begin{itemize}
  \item Probably binds to and stabilises inactivated sodium channels.
  \item Prevents further generation of action potentials involved in seizure activity.
  \item May also:
  \begin{itemize}
    \item reduce calcium entry into neurones
    \item block transmitter release
    \item enhance GABA action
  \end{itemize}
\end{itemize}

\subsubsection{Effects}
\begin{itemize}
  \item \textbf{Idiosyncratic}
  \begin{itemize}
    \item acne
    \item coarsening of facial features
    \item hirsutism
    \item gum hyperplasia
    \item folate-dependent megaloblastic anaemia
    \item aplastic anaemia
    \item skin rashes
    \item peripheral neuropathy
  \end{itemize}
  \item \textbf{Dose-related}
  \begin{itemize}
    \item ataxia
    \item nystagmus
    \item paraesthesia
    \item vertigo
    \item slurred speech
    \item rapid undiluted IV administration may cause hypotension and heart block
  \end{itemize}
  \item \textbf{Teratogenic}
  \begin{itemize}
    \item craniofacial abnormalities
    \item growth retardation
    \item limb defects
    \item cardiac defects
    \item mental retardation
  \end{itemize}
  \item \textbf{Drug interactions}
  \begin{itemize}
    \item induces hepatic mixed function oxidases
    \item increases metabolism of:
    \begin{itemize}
      \item warfarin
      \item benzodiazepines
      \item oral contraceptive pill
    \end{itemize}
    \item metabolism inhibited by:
    \begin{itemize}
      \item metronidazole
      \item chloramphenicol
      \item isoniazid
    \end{itemize}
    \item can then rise to toxic levels
    \item metabolism induced by:
    \begin{itemize}
      \item carbamazepine
      \item alcohol
    \end{itemize}
    \item then plasma levels may fall
  \end{itemize}
\end{itemize}

\subsubsection{Kinetics}
\begin{itemize}
  \item Oral bioavailability about 90\%.
  \item Protein binding about 90\%.
  \item Saturable hepatic hydroxylation.
  \item Shows zero-order kinetics just above the therapeutic range.
  \item Can induce its own metabolism and that of other drugs.
  \item Major metabolite excreted in urine.
\end{itemize}

\subsection{Carbamazepine}

\subsubsection{Uses}
\begin{itemize}
  \item Also used in trigeminal neuralgia.
  \item Mode of action similar to phenytoin.
  \item Can be given orally or rectally.
\end{itemize}

\subsubsection{Effects}
\begin{itemize}
  \item \textbf{CNS}
  \begin{itemize}
    \item headache
    \item diplopia
    \item ataxia
    \item vomiting
    \item drowsiness
    \item these often limit use
  \end{itemize}
  \item \textbf{Metabolic}
  \begin{itemize}
    \item antidiuretic effect leading to water retention
  \end{itemize}
  \item \textbf{Miscellaneous}
  \begin{itemize}
    \item drug-induced hepatitis
    \item rash in 5--10\%
    \item rarely agranulocytosis
  \end{itemize}
  \item \textbf{Teratogenic}
  \begin{itemize}
    \item facial abnormalities
    \item IUGR
    \item microcephaly
    \item mental retardation
    \item incidence increases with dose
    \item overall incidence stated as 1/300--2000 live births
  \end{itemize}
  \item \textbf{Drug interactions}
  \begin{itemize}
    \item enzyme induction gives many interactions like phenytoin
    \item concurrent phenytoin levels may rise or fall
    \item erythromycin can increase carbamazepine levels
  \end{itemize}
\end{itemize}

\subsubsection{Kinetics}
\begin{itemize}
  \item Well absorbed from gut.
  \item High oral bioavailability.
  \item About 75\% protein-bound.
  \item Extensive hepatic metabolism to carbamazepine 10,11-epoxide.
  \item This metabolite retains about 30\% of anticonvulsant activity.
  \item Strong hepatic enzyme inducer.
  \item Induces its own metabolism.
  \item Excreted almost entirely in urine as unconjugated metabolites.
\end{itemize}

\subsection{Sodium valproate}

\subsubsection{Uses}
\begin{itemize}
  \item Various forms of epilepsy including absence seizures.
  \item Also used in trigeminal neuralgia.
\end{itemize}

\subsubsection{Mechanism of action}
\begin{itemize}
  \item Appears to:
  \begin{itemize}
    \item stabilise inactive sodium channels
    \item stimulate central GABA-ergic inhibitory pathways
  \end{itemize}
  \item Generally well tolerated.
\end{itemize}

\subsubsection{Effects}
\begin{itemize}
  \item \textbf{Abdominal}
  \begin{itemize}
    \item nausea
    \item gastric irritation
    \item pancreatitis
    \item potentially fatal hepatotoxicity
  \end{itemize}
  \item \textbf{Haematological}
  \begin{itemize}
    \item thrombocytopenia
    \item reduced platelet aggregation
  \end{itemize}
  \item \textbf{Miscellaneous}
  \begin{itemize}
    \item transient hair loss
  \end{itemize}
  \item \textbf{Teratogenic}
  \begin{itemize}
    \item neural tube defects
  \end{itemize}
\end{itemize}

\subsubsection{Kinetics}
\begin{itemize}
  \item Well absorbed orally.
  \item Highly protein-bound ($\sim$90\%).
  \item Hepatic metabolism.
  \item Some metabolites are active.
  \item Excreted in urine.
\end{itemize}

\subsection{Phenobarbital}
\begin{itemize}
  \item Effective anticonvulsant.
  \item Use limited by significant sedation.
  \item Long-acting barbiturate.
  \item Induces hepatic enzymes.
  \item Interacts with:
  \begin{itemize}
    \item warfarin
    \item oral contraceptives
    \item other anticonvulsants
  \end{itemize}
\end{itemize}

\subsection{Benzodiazepines}
\begin{itemize}
  \item Referred back to earlier benzodiazepine section.
  \item Widely used in emergency treatment of status epilepticus.
  \item Act by enhancing the chloride-gating function of GABA.
\end{itemize}

\subsection{Lamotrigine}

\subsubsection{Uses}
\begin{itemize}
  \item Adjunct therapy.
  \item Used in a range of seizure types.
  \item In adults and in children down to 2 years.
  \item Also used in bipolar disorder.
  \item Also used as a mood stabiliser.
\end{itemize}

\subsubsection{Mechanism of action}
\begin{itemize}
  \item Chemically unrelated to other anticonvulsants.
  \item Works by inactivation of voltage-sensitive sodium channels.
  \item Specifically in neurones that produce excitatory neurotransmitters:
  \begin{itemize}
    \item glutamate
    \item aspartate
  \end{itemize}
  \item The chapter notes it may have additional actions not seen with classic sodium-channel blockers like phenytoin.
\end{itemize}

\subsubsection{Effects}
\begin{itemize}
  \item \textbf{Abdominal}
  \begin{itemize}
    \item hepatotoxicity
    \item pancreatitis
    \item nausea
  \end{itemize}
  \item \textbf{Dermatological}
  \begin{itemize}
    \item associated with life-threatening skin reactions
  \end{itemize}
  \item \textbf{Neurological}
  \begin{itemize}
    \item rare reports of neuroleptic malignant syndrome
  \end{itemize}
\end{itemize}

\subsubsection{Kinetics}
\begin{itemize}
  \item Half-life 14 h.
  \item Oral bioavailability 98\%.
  \item Conjugated to an inactive metabolite.
\end{itemize}

\subsection{Gabapentin}

\subsubsection{Uses}
\begin{itemize}
  \item Focal seizures:
  \begin{itemize}
    \item 900--1200 mg/day
  \end{itemize}
  \item Neuropathic pain:
  \begin{itemize}
    \item 1800 mg/day
  \end{itemize}
  \item Dose should be titrated up and down at start and end of therapy.
  \item Has been used as premedication in an attempt to reduce postoperative pain.
\end{itemize}

\subsubsection{Mechanism of action}
\begin{itemize}
  \item Acts on the presynaptic $\alpha_2\delta$ subunit of voltage-gated calcium channels in cortical neurones.
  \item Increases synaptic concentration of GABA.
  \item May also have some NMDA antagonist activity.
\end{itemize}

\subsubsection{Effects}
\begin{itemize}
  \item \textbf{Infective}
  \begin{itemize}
    \item viral infections common
  \end{itemize}
  \item \textbf{Neurological}
  \begin{itemize}
    \item psychiatric symptoms
    \item ataxia
    \item dizziness
    \item fatigue
    \item fever
  \end{itemize}
\end{itemize}

\subsubsection{Kinetics}
\begin{itemize}
  \item Oral bioavailability about 60\%.
  \item Bioavailability decreases with increasing dose.
  \item Not protein-bound.
  \item Not metabolised.
  \item Easier to combine with other anticonvulsants for this reason.
  \item Eliminated solely by the kidneys.
  \item Half-life 6 h.
\end{itemize}

\subsection{Vigabatrin}

\subsubsection{Uses}
\begin{itemize}
  \item Adjunct with other antiepileptics for complex partial seizures.
  \item Maximum dose 1.5 g bd.
\end{itemize}

\subsubsection{Mechanism of action}
\begin{itemize}
  \item Irreversible inhibition of GABA-transaminase.
  \item Therefore reduces GABA breakdown.
  \item Duration of action 24 h.
\end{itemize}

\subsubsection{Kinetics}
\begin{itemize}
  \item Excreted unchanged in urine.
  \item Interacts with phenytoin by an unknown mechanism.
  \item Can reduce phenytoin concentration by up to 25\%.
  \item Elimination half-life 6 h.
\end{itemize}

\subsection{Pregabalin}

\subsubsection{Uses}
\begin{itemize}
  \item Adjunct for partial seizures in adults.
  \item Also used in anxiety.
  \item Described as a specifically designed successor to gabapentin.
\end{itemize}

\subsubsection{Mechanism of action}
\begin{itemize}
  \item Binds the $\alpha_2\delta$ subunit of voltage-gated calcium channels.
  \item Similar to gabapentin.
  \item Unlike gabapentin, it does not appear to alter synaptic GABA concentration.
\end{itemize}

\subsubsection{Effects}
\begin{itemize}
  \item Central effects common:
  \begin{itemize}
    \item dizziness
    \item drowsiness
  \end{itemize}
  \item Other unpleasant effects listed:
  \begin{itemize}
    \item weight gain
    \item confusion
    \item ataxia
    \item GI upset
    \item erectile dysfunction
  \end{itemize}
\end{itemize}

\subsubsection{Kinetics}
\begin{itemize}
  \item Rapidly absorbed.
  \item Oral bioavailability 90\%.
  \item Unaffected by dose.
  \item Absorption delayed by food.
  \item Not protein-bound.
  \item Not metabolised.
  \item Eliminated in urine.
\end{itemize}

\end{document}
